% Physical unit 018; accepted Session 02.
\GordanUnitStart{018}{printed page 13}
\GordanTopFolio{13}
Integrale
\[
 \int_{\psi_1}^{\psi_2}d\psi\{(11)\delta l'\delta l'+2(01)\delta l\delta l'+(00)\delta l\delta l\}
\]
tamdiu igitur positivum est, quoad constantes $\alpha$ et $\beta$ hac ratione possint determinari, ut $N$ in spatio a $\psi_1$ ad $\psi_2$ nunquam infinitum evadat.

\begin{center}\S\ 7.\end{center}

Quocirca functionem $N$ transformare placet.

Quam ad rem differentiamus aequationem lineae geodeticae:
\[
 \psi-\psi_1=\int_{\varphi_1}^{\varphi}{d\varphi\cos l_0\sqrt{1-e^2+e^2\sin^2l_0\cos^2\varphi}\over 1-\sin^2l_0\cos^2\varphi}
\]
(in qua angulos $\varphi_1$ et $\varphi$ constat definiri sic:
\[
 \sin l_0\cos\varphi_1=\sin l,
 \qquad \sin l_0\cos\varphi=\sin l)
\]
secundum $\psi_1$ et $l_0$ partialiter, qua in differentiatione $\psi$ constans sumendum est.

Unde sequuntur identitates:
\[
 -1={\cos l_0\sqrt{1-e^2+e^2\sin^2l_0\cos^2\varphi}\over 1-\sin^2l_0\cos^2\varphi}{d\varphi\over d\psi_1}
 =-\cot l_0{\sqrt{1-e^2\cos^2l}\over \cos l\sin\varphi}{dl\over d\psi_1}
\]
\[
\begin{aligned}
0={}&e^2\sin^2l_0\int_{\varphi_1}^{\varphi}d\varphi\,{\cos^2\varphi\over \sqrt{1-e^2+e^2\sin^2l_0\cos^2\varphi}}
       +\cot\varphi\sqrt{1-e^2\cos^2l}-\cot\varphi_1\sqrt{1-e^2\cos^2l_1} \\
&{}-\cos l_0{\sqrt{1-e^2\cos^2l}\over \cos l\sin\varphi}\cdot {dl\over dl_0},
\end{aligned}
\]
ex quibus posito:
\[
\begin{aligned}
&e^2\sin^2l_0\sin\varphi\int_{\varphi_1}^{\varphi}d\varphi\,{\cos^2\varphi\over \sqrt{1-e^2+e^2\sin^2l_0\cos^2\varphi}}\\
&{}+\cos\varphi\sqrt{1-e^2\cos^2l}-\cot\varphi_1\sin\varphi\sqrt{1-e^2\cos^2l_1}=u.
\end{aligned}
\]
\[
 {\cos l_0\over \cos l}\sqrt{1-e^2\cos^2l}=v
\]
quotientes differentialium inveniuntur:
\[
 {dl\over d\psi}={\sin l_0\sin\varphi\over v};\qquad {dl\over dl_0}={u\over v}
\]
quibus valoribus ad $N$ relatis, fit:
\[
\begin{aligned}
-N&=(01)+(11){\left({\alpha{\sin l_0\sin\varphi\over v}+\beta{u\over v}}\right)'
\over {\alpha{\sin l_0\sin\varphi\over v}+\beta{u\over v}}}\\
&=(01)+{(11)\cdot(1-\sin^2l_0\cos^2\varphi)\over \cos l_0\sqrt{1-e^2+e^2\sin^2l_0\cos^2\varphi}}
\cdot {v{d\over d\varphi}\left\{\alpha{\sin l_0\sin\varphi\over v}+\beta{u\over v}\right\}
\over \alpha\sin l_0\sin\varphi+\beta u}.
\end{aligned}
\]
